\documentclass[10pt,a4paper]{extarticle}
\usepackage{parskip}

\usepackage[utf8]{inputenc}
\usepackage[T1]{fontenc}
\usepackage[brazil]{babel}
\usepackage{amsmath, amssymb}
\usepackage{geometry}
\usepackage{xcolor}
\usepackage{hyperref}
\usepackage{enumitem}
\usepackage{microtype}
\usepackage{csquotes}
\usepackage{natbib}

\definecolor{important}{RGB}{255, 100, 100}
\definecolor{notice}{RGB}{0, 102, 204}

\newcommand{\impt}[1]{\textcolor{important}{#1}}
\newcommand{\ntc}[1]{\textcolor{notice}{#1}}

\setlist{nosep}
\geometry{margin=2cm}

\hypersetup{
  colorlinks=true,
  linkcolor=notice,
  citecolor=notice,
  urlcolor=notice
}

\title{Estudo Dirigido II -- Falsificacionismo\footnote{Questões baseadas em \citep{chalmers1993}.}}
\author{BCC502 -- Metodologia Científica para Ciência da Computação}
\date{\today}

\begin{document}

\vspace{-1cm}
\maketitle

\section*{Capítulo 4}
\begin{enumerate}[label=\textbf{\arabic*.}, itemsep=6pt]
  \item Qual é a visão do falsificacionista a respeito das observações e teorias?
  \item Qual é a visão do falsificacionista em relação ao progresso da ciência?
  \item Como o falsificacionismo pode ser justificado logicamente?
  \item O que é uma hipótese falsificável?
  \item Como se estabelece a preferência entre uma teoria e outra dentro do falsificacionismo?
  \item Qual é a consequência atrativa que vem da exigência de que as teorias sejam altamente falsificáveis?
  \item Como se dá o progresso da ciência na visão do falsificacionista?
\end{enumerate}

\section*{Capítulo 5}
\begin{enumerate}[resume, label=\textbf{\arabic*.}, itemsep=6pt]
  \item Quais condições uma hipótese deve satisfazer?
  \item Qual é a explicação falsificacionista sofisticada da ciência?
  \item Qual é o problema técnico do falsificacionismo ingênuo? Como o falsificacionismo sofisticado o resolve?
  \item O que é uma modificação \emph{ad hoc}? Apresente um exemplo de uma modificação \emph{ad hoc} e um de uma modificação não \emph{ad hoc}.
  \item De onde advêm as principais contribuições ao crescimento do conhecimento científico? Por quê?
  \item Qual é a meta da ciência no falsificacionismo?
  \item Como se diferem as visões indutivistas e falsificacionistas sobre o papel da observação?
\end{enumerate}

\section*{Capítulo 6}
\begin{enumerate}[resume, label=\textbf{\arabic*.}, itemsep=6pt]
  \item Como se diferem o \enquote{falsificacionista ingênuo} e o \enquote{falsificacionista mais sofisticado}?
  \item Justifique a afirmação a seguir: \enquote{Todas as proposições são falíveis}.
  \item Quais são as objeções feitas ao falsificacionismo?
  \item Como Popper defende o falsificacionismo dessas objeções? Por que essa defesa não salva o falsificacionismo? Dê exemplos.
  \item Quando uma proposição de afirmação é aceitável?
  \item Por que a aplicação do falsificacionismo em situações de teste realistas é complexa? Dê exemplos.
  \item Quais são as evidências da inadequação histórica do falsificacionismo?
  \item Por que nem falsificacionismo nem indutivismo são compatíveis com a revolução copernicana?
\end{enumerate}

\bibliographystyle{apalike}
\bibliography{ref}
\end{document}